\clearpage
\phantomsection
\addcontentsline{toc}{chapter}{{Mở đầu}}
\chapter*{Mở đầu}

\section{Lý do chọn đề tài}

Trong ngành công nghiệp trò chơi điện tử hiện đại, việc cân bằng gameplay là một trong những thách thức lớn nhất mà các nhà phát triển phải đối mặt. Đặc biệt với thể loại Roguelike Deckbuilder, nơi mỗi lần chơi là một trải nghiệm độc đáo với các lựa chọn ngẫu nhiên về thẻ bài, di vật và kẻ thù, việc đảm bảo game vừa thách thức vừa công bằng là cực kỳ khó khăn \cite{schell2019art}.

Phương pháp cân bằng game truyền thống thường dựa vào kinh nghiệm và thử nghiệm thủ công của các nhà thiết kế. Tuy nhiên, với số lượng lớn các tham số cần điều chỉnh (giá trị sát thương của thẻ bài, máu của kẻ địch, hiệu ứng của di vật, v.v.), phương pháp này tốn rất nhiều thời gian và có thể bỏ sót các kịch bản edge case quan trọng \cite{jaffe2012optimizing}.

Xuất phát từ thực tiễn đó, đề tài này tập trung nghiên cứu và phát triển một hệ thống cân bằng game tự động sử dụng \textbf{Thuật toán di truyền (Genetic Algorithm - GA)} để tối ưu hóa các tham số game. Hệ thống không chỉ giúp tự động hóa quá trình cân bằng mà còn có khả năng khám phá các cấu hình tối ưu mà con người có thể không nghĩ tới \cite{holland1992adaptation, yannakakis2018artificial}.

\section{Mục tiêu nghiên cứu}

Mục tiêu chính của đề tài là xây dựng một hệ thống hoàn chỉnh bao gồm việc phát triển một trò chơi Roguelike Deckbuilder prototype với đầy đủ cơ chế gameplay cơ bản, xây dựng hệ thống mô phỏng tự động có khả năng chạy hàng nghìn lượt chơi để thu thập dữ liệu, và thiết kế cùng triển khai thuật toán di truyền để tối ưu hóa các tham số game.

Bên cạnh đó, hệ thống cần so sánh hiệu quả của hai phương pháp: \textbf{GA đơn mục tiêu} (tối ưu trực tiếp từng tham số bằng GA truyền thống) và \textbf{NSGA-II đa mục tiêu} (sử dụng thuật toán NSGA-II để tối ưu đa mục tiêu với không gian tham số được cấu trúc hóa) \cite{deb2002fast, zitzler2000comparison}. Cuối cùng, đề tài phát triển công cụ trực quan hóa để phân tích kết quả và đánh giá chất lượng cân bằng đạt được.

\section{Phương pháp nghiên cứu}

Để giải quyết bài toán phức tạp nêu trên, đề tài áp dụng cách tiếp cận đa chiều kết hợp giữa nghiên cứu lý thuyết, phát triển hệ thống và tối ưu hóa thực nghiệm.

\subsection{Nghiên cứu lý thuyết}
Giai đoạn nghiên cứu lý thuyết tập trung khảo sát các thuật toán di truyền (GA) và các biến thể hiện đại như NSGA-II \cite{deb2002fast, zitzler2000comparison}, nghiên cứu các phương pháp đánh giá chất lượng cân bằng game thông qua các chỉ số như tỷ lệ thắng, điểm HP trung bình khi chiến thắng, và tầng trung bình khi thua. Đồng thời, đề tài cũng tìm hiểu về game design patterns trong thể loại Roguelike Deckbuilder, lấy cảm hứng từ các tựa game thành công như \textit{Slay the Spire} \cite{slaythespire}.

\subsection{Phát triển hệ thống}
Hệ thống được xây dựng với game engine bao gồm các module: Combat System, Deck Management, Map Generation, và Room Handlers. AI agent được thiết kế để tự động chơi game (HeuristicPlayerAI), cùng với hệ thống mô phỏng có khả năng chạy song song và ghi log đầy đủ \cite{yannakakis2018artificial}.

\subsection{Tối ưu hóa và thực nghiệm}
Quá trình tối ưu hóa được triển khai theo hai hướng tiếp cận: GA đơn mục tiêu với GA truyền thống xử lý khoảng 500 tham số, và NSGA-II đa mục tiêu chỉ cần khoảng 50 tham số phân cấp. Cả hai phương pháp đều được chạy thực nghiệm và so sánh kết quả, sau đó phân tích Pareto front để tìm các giải pháp cân bằng tối ưu \cite{preuss2012multiobjective}.

\section{Nội dung nghiên cứu}

Dự án này tập trung xây dựng một Roguelike Deckbuilder game hoàn chỉnh với các cơ chế chiến đấu theo lượt (turn-based combat), quản lý bộ bài (deck building), di vật (relics), và bản đồ phòng (room-based map) \cite{dormans2011level}. Hệ thống AI tự động chơi game được phát triển dựa trên heuristics để có thể mô phỏng hàng nghìn lượt chơi, thu thập dữ liệu về gameplay metrics.

Không gian tham số tối ưu hóa bao gồm sát thương và giá mana của thẻ bài, giá trị sức khỏe và sát thương của kẻ địch, cùng với hiệu ứng của di vật và status effects. Hàm fitness đa mục tiêu được xây dựng để đánh giá chất lượng cân bằng theo ba tiêu chí: \textbf{Balance} (độ cân bằng giữa thắng/thua), \textbf{Engagement} (độ hấp dẫn của gameplay), và \textbf{Coherence} (tính nhất quán của thiết kế) \cite{preuss2012multiobjective}.

Hai phương pháp tối ưu hóa với các chiến lược khác nhau được triển khai và so sánh, cùng với công cụ trực quan hóa để phân tích kết quả including visualizer cho Pareto front 3D và evolution charts.

\section{Đóng góp của đề tài}

Về mặt khoa học, đề tài đề xuất một mô hình tham chiếu cho việc ứng dụng thuật toán tiến hóa vào bài toán cân bằng game, đặc biệt là thể loại Roguelike Deckbuilder \cite{togelius2011search}. Nghiên cứu so sánh giữa phương pháp tối ưu trực tiếp và tối ưu phân cấp trong bối cảnh game balancing, cung cấp insights về hiệu quả của hierarchical parameter grouping.

Về mặt thực tiễn, đề tài cung cấp một công cụ hoàn chỉnh giúp các nhà phát triển game có thể tự động hóa quá trình cân bằng, tiết kiệm hàng trăm giờ thử nghiệm thủ công \cite{jaffe2012optimizing, nelson2016towards}. Hệ thống có khả năng mở rộng và có thể áp dụng cho các loại game khác với các điều chỉnh phù hợp.

Về mặt công nghệ, đề tài kiểm chứng khả năng kết hợp giữa game development và evolutionary algorithms trong một hệ thống thực tế. Công cụ trực quan hóa được phát triển giúp phân tích và hiểu sâu về quá trình tiến hóa của các tham số game, hỗ trợ game designers trong việc đưa ra quyết định thiết kế dựa trên dữ liệu khách quan \cite{smith2011analyzing}.

\section{Bố cục của báo cáo}
Báo cáo được tổ chức thành các chương chính sau:
\begin{itemize}
    \item \textbf{Mở đầu}: Trình bày bối cảnh, lý do chọn đề tài, mục tiêu và phương pháp tiếp cận.
    \item \textbf{Chương 1 - Tổng quan và Cơ sở lý thuyết}: Giới thiệu về thể loại Roguelike Deckbuilder, khảo sát các thuật toán di truyền (GA, NSGA-II), và phân tích các nghiên cứu liên quan về game balancing.
    \item \textbf{Chương 2 - Phân tích và Thiết kế hệ thống}: Mô tả kiến trúc chi tiết của game engine, hệ thống mô phỏng, và các module tối ưu hóa.
    \item \textbf{Chương 3 - Cài đặt và Triển khai}: Trình bày quy trình hiện thực hóa các module, các thách thức kỹ thuật gặp phải và giải pháp xử lý.
    \item \textbf{Chương 4 - Thực nghiệm và Đánh giá}: Phân tích kết quả chạy thực nghiệm với cả hai phương pháp tối ưu hóa, so sánh hiệu quả và đánh giá chất lượng cân bằng đạt được.
    \item \textbf{Chương 5 - Kết luận và Hướng phát triển}: Tổng kết các kết quả đạt được, hạn chế của hệ thống và đề xuất hướng mở rộng trong tương lai.
\end{itemize}